% P11 paper module: R24 uniform coercivity and resolvent bridge.

\subsection{Uniform coercivity and the inverse-root resolvent bridge}
\label{subsec:o3w-resolvent-bridge}

The R22 metric representation contains the nested inverse square roots
$A_R^{-1/2}$ and $A_S^{-1/2}$.  For the concrete P11 relative metrics their norms are,
however, uniformly controlled at fixed base terminal.  This separates the lower-spectral
question from the remaining orientation problem and gives an exact resolvent formula for
the latter.

\begin{proposition}[Uniform coercivity of the P11 relative metrics]
\label{prop:o3w-coercivity}
Fix $0<R<S<T_0$ and let $X\in\{R,S\}$.  For every $U>T_0$, put
\[
 A_X:=A_{T_0,U}^{X,-}.
\]
Then
\[
 \boxed{
 a_0 I\le A_X,
 \qquad
 a_0:=\frac1{1+\|H_{T_0}\|^2}>0.
 }
\tag{O3W.1}\label{eq:o3w-lower}
\]
Consequently
\[
 \boxed{
 \|A_X^{-1/2}\|
 \le a_0^{-1/2}
 =\sqrt{1+\|H_{T_0}\|^2}
 }
\tag{O3W.2}\label{eq:o3w-invroot-bound}
\]
uniformly in $U$.
Together with Lemma~\ref{lem:o3l-relative-upper},
\[
 a_0I\le A_X\le(1+\|H_U\|^2)I
 \ll\frac{e^U}{U}I.
\tag{O3W.3}\label{eq:o3w-two-sided}
\]
Thus the growth of the relative condition number comes from the upper spectral edge,
not from collapse of the lower edge.
\end{proposition}

\begin{proof}
For every nonzero source vector $z$, the Rayleigh quotient of $A_X$ is
\[
 \frac{q_U^X(E_{X,U}z)}{q_{T_0}^X(E_{X,T_0}z)}.
\]
Exact Gamma compatibility and positivity of the Schur term give
\[
 q_U^X(E_{X,U}z)\ge\mathfrak c_{\Gamma,X}[z].
\]
At the fixed terminal $T_0$, the graph-norm equivalence
\eqref{eq:graph-norm-equiv} gives
\[
 q_{T_0}^X(E_{X,T_0}z)
 \le(1+\|H_{T_0}\|^2)\mathfrak c_{\Gamma,X}[z].
\]
Their quotient is therefore at least $a_0$, proving~\eqref{eq:o3w-lower}.
Equation~\eqref{eq:o3w-invroot-bound} follows by functional calculus, and the upper
bound in~\eqref{eq:o3w-two-sided} is Lemma~\ref{lem:o3l-relative-upper}.
\end{proof}

\begin{remark}[Concrete inverse-root firewall]
\label{rem:o3w-concrete-firewall}
Proposition~\ref{prop:o3n-rankone-insufficiency} remains a valid abstract information
firewall for entrywise rank-one leading data, but the concrete P11 family has the extra
constraint~\eqref{eq:o3w-lower}.  Hence one must not interpret the remaining P11 gate as
possible blow-up of $\|A_X^{-1/2}\|$ caused by
$\lambda_{\min}(A_X)\downarrow0$.  The inverse-root norms are uniformly bounded.  What
remains unresolved is their orientation across the nested $R/S$ levels.
\end{remark}

Keep
\[
 A_R:=A_{T_0,U}^{R,-},
 \qquad
 A_S:=A_{T_0,U}^{S,-},
 \qquad
 W:=W_{R,S,-}^{[T_0]},
 \qquad
 P:=WW^*.
\]
The exact relative compression identity gives
\[
 A_R=W^*A_SW.
\]
Therefore
\[
 \boxed{
 \mathscr B_U
 :=(I-P)A_SW
 =A_SW-WA_R.
 }
\tag{O3W.4}\label{eq:o3w-B-intertwining}
\]
This is the same second-moment off-block used in O3l and O3r.

\begin{proposition}[Inverse-square-root resolvent bridge]
\label{prop:o3w-inverse-resolvent}
One has the exact operator-norm Bochner integral
\[
 \boxed{
 \begin{aligned}
 A_S^{-1/2}W-WA_R^{-1/2}
 ={}&-\frac1\pi\int_0^\infty t^{-1/2}
 (A_S+tI)^{-1}\mathscr B_U(A_R+tI)^{-1}\,dt.
 \end{aligned}
 }
\tag{O3W.5}\label{eq:o3w-inverse-resolvent}
\]
In particular
\[
 \boxed{
 \|A_S^{-1/2}W-WA_R^{-1/2}\|
 \le\frac1{2a_0^{3/2}}\|\mathscr B_U\|.
 }
\tag{O3W.6}\label{eq:o3w-inverse-bound}
\]
\end{proposition}

\begin{proof}
For a positive invertible operator $A$,
\[
 A^{-1/2}
 =\frac1\pi\int_0^\infty t^{-1/2}(A+tI)^{-1}\,dt.
\]
The resolvent quasi-commutator satisfies
\[
\begin{aligned}
 &(A_S+tI)^{-1}W-W(A_R+tI)^{-1}\\
 &\qquad
 =(A_S+tI)^{-1}
 \bigl(W(A_R+tI)-(A_S+tI)W\bigr)
 (A_R+tI)^{-1}\\
 &\qquad
 =-(A_S+tI)^{-1}\mathscr B_U(A_R+tI)^{-1},
\end{aligned}
\]
which proves~\eqref{eq:o3w-inverse-resolvent}.  By
Proposition~\ref{prop:o3w-coercivity},
\[
 \|(A_X+tI)^{-1}\|\le(a_0+t)^{-1}.
\]
Hence the integral converges in norm, and
\[
 \frac1\pi\int_0^\infty
 \frac{t^{-1/2}}{(a_0+t)^2}\,dt
 =\frac1{2a_0^{3/2}},
\]
which gives~\eqref{eq:o3w-inverse-bound}.
\end{proof}

\begin{proposition}[Square-root resolvent bridge]
\label{prop:o3w-square-resolvent}
The corresponding positive-square-root intertwining defect is
\[
 \boxed{
 \begin{aligned}
 A_S^{1/2}W-WA_R^{1/2}
 ={}&\frac1\pi\int_0^\infty t^{1/2}
 (A_S+tI)^{-1}\mathscr B_U(A_R+tI)^{-1}\,dt.
 \end{aligned}
 }
\tag{O3W.7}\label{eq:o3w-square-resolvent}
\]
Consequently
\[
 \boxed{
 \|A_S^{1/2}W-WA_R^{1/2}\|
 \le\frac1{2\sqrt{a_0}}\|\mathscr B_U\|.
 }
\tag{O3W.8}\label{eq:o3w-square-bound}
\]
\end{proposition}

\begin{proof}
Use
\[
 A^{1/2}
 =\frac1\pi\int_0^\infty
 t^{-1/2}A(A+tI)^{-1}\,dt
\]
and
\[
 A(A+tI)^{-1}=I-t(A+tI)^{-1}.
\]
Thus
\[
\begin{aligned}
 &A_S(A_S+tI)^{-1}W
 -WA_R(A_R+tI)^{-1}\\
 &\qquad
 =-t\bigl((A_S+tI)^{-1}W-W(A_R+tI)^{-1}\bigr)\\
 &\qquad
 =t(A_S+tI)^{-1}\mathscr B_U(A_R+tI)^{-1}.
\end{aligned}
\]
Integration gives~\eqref{eq:o3w-square-resolvent}.  Finally
\[
 \frac1\pi\int_0^\infty
 \frac{t^{1/2}}{(a_0+t)^2}\,dt
 =\frac1{2\sqrt{a_0}},
\]
proving~\eqref{eq:o3w-square-bound}.
\end{proof}

\begin{corollary}[Common resolvent source of the positive and inverse-root off-blocks]
\label{cor:o3w-common-source}
Applying $I-P$ to~\eqref{eq:o3w-square-resolvent} gives
\[
 \boxed{
 \mathscr C_U
 =(I-P)A_S^{1/2}W
 =\frac1\pi\int_0^\infty t^{1/2}
 (I-P)(A_S+tI)^{-1}\mathscr B_U(A_R+tI)^{-1}\,dt.
 }
\tag{O3W.9}\label{eq:o3w-C-resolvent}
\]
Likewise
\[
 \boxed{
 (I-P)A_S^{-1/2}W
 =-\frac1\pi\int_0^\infty t^{-1/2}
 (I-P)(A_S+tI)^{-1}\mathscr B_U(A_R+tI)^{-1}\,dt.
 }
\tag{O3W.10}\label{eq:o3w-inverse-offblock}
\]
Hence O3r's square-root witness and the inverse-root orientation entering O3t/O3u are
different spectral weightings of the same second-moment off-block.
\end{corollary}

\begin{remark}[R24 reduction]
\label{rem:o3w-reduction}
The raw norm lower bound for $\mathscr B_U$ from O3l does not decide the inverse-root
integral: the two resolvents may strongly suppress precisely the large directions.  But
R24 replaces the vague phrase ``inverse-square-root orientation'' by the explicit
resolvent-smoothed quantity
\[
 (A_S+tI)^{-1}\mathscr B_U(A_R+tI)^{-1}.
\]
This is a concrete analytic object built from operators already present in P11; no new
auxiliary gauge is introduced.
\end{remark}

\begin{openproblem}[R24-F: fixed-vector resolvent-weighted off-block asymptotics]
\label{open:o3w-resolvent-offblock}
For fixed smooth odd source vectors $f$, determine the asymptotics of
\[
 \int_0^\infty t^{-1/2}
 (A_S+tI)^{-1}\mathscr B_U(A_R+tI)^{-1}f\,dt
\]
and of the corresponding $t^{1/2}$-weighted integral as $U\to\infty$.
A vanishing result on a fixed dense core would give a genuine nested inverse-root
compatibility statement; a fixed-vector positive limsup would give an inverse-root
orientation obstruction.  The existing raw second-moment witness alone decides neither
alternative.
\end{openproblem}
